%! TeX program = lualatex
% The two facts in the limits-at-infinity box of lessons/limit-at-infinity.tex.
% Each page carries a fast and a slow power, so the picture reads as "every
% positive power behaves this way" rather than "this one curve does", plus the
% exponential, which outruns (or outdecays) every power.
%
% Only x > 0 is drawn: the exponent is any positive constant, and x^{1/2}
% and friends are undefined to the left of the origin.
\documentclass[tikz]{standalone}
\input{../typography.tex.preamble}
\input{../tikz.tex.preamble}

\pgfplotsset{fact plot/.style={
    basic curve,
    % equal scaling: one unit is the same length on both axes, so the shapes
    % of the curves are comparable between the two pictures
    axis equal image,
    width=1.9in,
    xmin=0, xmax=4.6,
    ymin=0, ymax=4.3,
    enlargelimits=false,
    xtick distance=2, ytick distance=2,
    clip=true,
    legend cell align=left,
    legend style={draw=none, fill=none, font=\scriptsize, inner sep=1pt,
                  row sep=-2pt},
}}

\begin{document}
% page 1 -- x^{positive constant} grows without bound
\begin{tikzpicture}
  \begin{axis}[fact plot, legend pos=south east]
    \addplot[very thick, attn, domain=0:1.47, samples=100] {exp(x)};
    \addplot[very thick, main, domain=0:2.07, samples=100] {x^2};
    \addplot[very thick, supp, domain=0:4.5, samples=100] {sqrt(x)};
    \legend{\(e^{x}\), \(x^{2}\), \(\sqrt{x}\)}
  \end{axis}
\end{tikzpicture}

% page 2 -- 1/x^{positive constant} decays to the horizontal asymptote y = 0
\begin{tikzpicture}
  \begin{axis}[fact plot, legend pos=north east]
    \addplot[thin, dashed, forget plot] coordinates {(0,0) (4.6,0)};
    \addplot[very thick, attn, domain=0:4.5, samples=100] {exp(-x)};
    \addplot[very thick, main, domain=0.485:4.5, samples=100] {1/x^2};
    \addplot[very thick, supp, domain=0.055:4.5, samples=100] {1/sqrt(x)};
    \legend{\(e^{-x}\), \(\frac{1}{x^{2}}\), \(\frac{1}{\sqrt{x}}\)}
  \end{axis}
\end{tikzpicture}
\end{document}
